In this section, we prove our first main result, which is as follows.

\begin{thm} \label{thm:rank-in-NC}
The problem of computing the (real) rank of an unambiguous matrix monoid is in~$\NC^2$.
\end{thm}

% Fix indentation
\noindent
In order to use \autoref{thm:pseud-linsyst}, we need the following lemma. We use the notation defined in the previous section.
Recall that we defined $U := \spn{\alpha^T M(w) - \alpha^T \mid w \in \Sigma^*}$.

\begin{lem} \label{lem:U-in-NC}
If $M$ is strongly connected,
one can compute a basis of~$U$ in~\NCT.
\end{lem}

% Fix indentation
\noindent
For each $a \in \Sigma$ define $M'(a) := M(a) - I \in \{-1,0,1\}^{Q \times Q}$ and extend~$M'$ to $M' : \Sigma^* \to \mathbb{Z}^{Q \times Q}$ by defining $M'(a_1 \cdots a_k) = M'(a_1) \cdots M'(a_k)$.
Define \[U' := \spn{\alpha^T M'(w) \mid w \in \Sigma^+}.\]
Note that here $w$ ranges over~$\Sigma^+$, i.e., nonempty words, only. By definition, $U'$ is closed under right multiplication by $M'(a)$ for all $a \in \Sigma$.
We first show the following lemma.

\begin{lem} \label{lem:U=U'}
We have $U = U'$.
\end{lem}
\begin{proof}
For the inclusion $U \subseteq U'$, we prove by induction on $i \ge 0$ that for all length-$i$ words $w \in \Sigma^i$ we have $\alpha^T (M(w) - I) \in U'$.
Concerning the induction base, $i=0$, we have $\alpha^T (M(\varepsilon) - I) = 0 \in U'$.
Concerning the induction step, let $i \ge 0$, and let $w \in \Sigma^i$ and $a \in \Sigma$.
We have
\begin{align*}
\alpha^T (M(w a) - I)
\ &= \ \alpha^T (M(w) - I) M(a) + \alpha^T (M(a) - I) \\
\ &= \ \alpha^T (M(w) - I) (M(a) - I) + \alpha^T (M(w) - I) + \alpha^T (M(a) - I) \\
\ &= \ \alpha^T (M(w) - I) M'(a) + \alpha^T (M(w) - I) + \alpha^T M'(a)\,.
\end{align*}
It holds that $\alpha^T M'(a) \in U'$, and, by the induction hypothesis, $\alpha^T (M(w) - I) \in U'$.
It follows that $\alpha^T (M(w a) - I) \in U'$.

For the converse, $U' \subseteq U$, we proceed similarly by induction.
Concerning the induction base, $i=1$, for all $a \in \Sigma$ we have $\alpha^T M'(a) = \alpha^T (M(a) - I) \in U$.
Concerning the induction step, let $i \ge 1$, and let $w \in \Sigma^i$ and $a \in \Sigma$.
By the induction hypothesis there are $n \le |Q|$ and $w_1, \ldots, w_n \in \Sigma^*$ and $\lambda_1, \ldots, \lambda_n \in \mathbb{R}$ such that $\alpha^T M'(w) = \sum_{i=1}^n \lambda_i \alpha^T (M(w_i) - I)$.
Thus, we have
\begin{align*}
\alpha^T M'(w a)
\ &= \ \alpha^T M'(w) (M(a) - I) 
\ = \ \sum_{i=1}^n \lambda_i \alpha^T (M(w_i) - I) (M(a) - I) \\
\ &= \ \sum_{i=1}^n \lambda_i \alpha^T \big((M(w_i a) - I) - (M(a) - I) - (M(w_i) - I) \big) \in \ U\,,
\end{align*}
as required.
\end{proof}


\begin{proof}[Proof of \autoref{lem:U-in-NC}]
For each $a \in \Sigma$ define $U'_a := \spn{\alpha^T M'(a) M'(w) \mid w \in \Sigma^*}$.
Using the technique from \cite[Section 4.2]{KieferMW14} (see \cite[Proposition 5.2]{KieferWeighted20} for a clearer explanation), for each $a \in \Sigma$ one can compute\footnote{In \cite{KieferMW14,KieferWeighted20} only membership in~$\NC$ is claimed, but the bottleneck computations are matrix powering and rank computation, which can in fact be done in $\mathrm{DET} \subseteq \NC^2$; see~\cite{Cook85}.
The computations in~\cite[Proposition 5.2]{KieferWeighted20} are on polynomially larger matrices, but this does not impact the membership in~$\NC^2$, as $\log^2(\mathit{poly}(n)) = O(\log^2(n))$.} a basis of~$U'_a$ in~$\NC^2$.
The union of these bases, say $\Gamma = \{\gamma_1^T, \ldots, \gamma_n^T\}$ for some $n \le |\Sigma| |Q|$, spans~$U'$, which equals~$U$ by \autoref{lem:U=U'}.
To shrink~$\Gamma$ to a basis of~$U$, for each $i \in \{1, \ldots, n\}$ include~$\gamma_i^T$ in the basis if and only if $\dim \spn{\gamma_1^T, \ldots, \gamma_{i-1}^T} < \dim \spn{\gamma_1^T, \ldots, \gamma_i^T}$.
The latter (rank) computation can be done in~$\NC^2$~\cite{IbarraMoranRosier80}. %; see also \cite[Section~5]{Cook85}.
\end{proof}

% Fix indentation
\noindent
Now we can prove \autoref{thm:rank-in-NC}.
\begin{proof}[Proof of \autoref{thm:rank-in-NC}]
Let $M : \Sigma \to \{0,1\}^{Q \times Q}$ be an unambiguous monoid morphism.
Its strongly connected components can be computed in $\NL \subseteq \NC^2$.
It follows from the proof of \cite[Proposition 3]{Kiefer2021} that one can check each component for completeness in~$\NC^2$, since a zero-one monoid contains the zero matrix if and only if the joint spectral radius of the set of its generators is strictly less than one \cite{Kiefer2021}.
Therefore, using \autoref{prop:unamb-sc}, we can assume in the rest of the proof that $M$ is complete and strongly connected.

We use the fact that one can compute a solution of a possibly singular linear system of equations in~\NCT \cite[Section~5]{BorodinGH82}.
First, compute in~\NCT vectors $\alpha, \beta \in \mathbb{Q}_{>0}^Q$ with $\alpha^T \overline{A} = \alpha^T$ and $\overline{A} \beta = \beta$ and $\alpha^T \beta = 1$.
Using \autoref{lem:U-in-NC} compute in~\NCT a basis of~$U$.
Choose an arbitrary $q \in Q$ and compute $\Mer(q)$ in $\NL \subseteq \NCT$ with a reachability analysis.
Then, solve the linear system from \autoref{thm:pseud-linsyst} to compute in~\NCT a maximal pseudo-column~$y \in \mathbb{Q}^Q$.
Hence, using \autoref{lem:Cas}, we can compute~$\mcw = \alpha^T y$ in~\NCT.
Symmetrically, we can compute~$\mrw$ in~\NCT.
Finally, by \autoref{thm:rank-weight}, we can compute $r = \frac{1}{\mcw \cdot \mrw}$ in~\NCT.
\end{proof}

\begin{exa} \label{ex:snap3}
We continue \autoref{ex:snap1} and \autoref{ex:snap2}.
Since $M$~is a total DFA, it is natural to take $\beta = [Q]$.
Note that $\alpha^T \beta = 1$.
Since every row is of the form $[q]^T$ for some~$q$, we have $\mrw = 1$.
Recall from \autoref{ex:snap2} that $\mcw = \frac12$.
With \autoref{thm:rank-weight} we conclude that $r = \frac{1}{\mcw \cdot \mrw} = 2$, as observed in \autoref{ex:snap1}.
\end{exa}
